\section{Classical Field Theory}

\subsection{Fields and the Action Principle}
A (classical) field is a function on spacetime:
\begin{equation}
\phi_a(x), \qquad x^\mu=(t,\mathbf{x}), \qquad a=1,\dots,N,
\end{equation}
where $a$ labels different field components (e.g. internal indices).

The dynamics is determined by an action functional
\begin{equation}
S[\phi] = \int d^4x\, \mathcal{L}\big(\phi_a,\partial_\mu\phi_a,x\big),
\end{equation}
where $\mathcal{L}$ is the \emph{Lagrangian density}. For local relativistic theories,
\begin{equation}
\mathcal{L}=\mathcal{L}\big(\phi_a,\partial_\mu\phi_a\big),
\end{equation}
i.e. no explicit $x$ dependence.

\subsection{Euler--Lagrange Equations for Fields}
Vary the action under $\phi_a\to \phi_a+\delta\phi_a$ with $\delta\phi_a$ vanishing on the boundary:
\begin{equation}
\delta S = 0 \quad \Rightarrow \quad
\frac{\partial\mathcal{L}}{\partial \phi_a}
-\partial_\mu\left(\frac{\partial\mathcal{L}}{\partial(\partial_\mu\phi_a)}\right)=0.
\end{equation}

\paragraph{Example: Klein--Gordon field.}
\begin{equation}
\mathcal{L}_{\mathrm{KG}}
=\frac{1}{2}\partial_\mu\phi\,\partial^\mu\phi-\frac{1}{2}m^2\phi^2.
\end{equation}
Then
\begin{equation}
(\square+m^2)\phi=0.
\end{equation}

\paragraph{Example: complex scalar.}
\begin{equation}
\mathcal{L}
=\partial_\mu\phi^\ast\,\partial^\mu\phi - m^2\phi^\ast\phi.
\end{equation}
Treat $\phi$ and $\phi^\ast$ as independent in variation:
\begin{equation}
(\square+m^2)\phi=0,\qquad (\square+m^2)\phi^\ast=0.
\end{equation}

\subsection{Canonical Momentum and Hamiltonian Density}
Define canonical conjugate momentum fields:
\begin{equation}
\pi_a(x) := \frac{\partial\mathcal{L}}{\partial(\partial_0\phi_a)}.
\end{equation}
Hamiltonian density:
\begin{equation}
\mathcal{H} := \pi_a \dot{\phi}_a - \mathcal{L},
\qquad \dot{\phi}_a:=\partial_0\phi_a.
\end{equation}
Total Hamiltonian:
\begin{equation}
H=\int d^3x\,\mathcal{H}.
\end{equation}

Hamilton's equations (field form):
\begin{equation}
\dot{\phi}_a(\mathbf{x},t)=\frac{\delta H}{\delta \pi_a(\mathbf{x},t)},
\qquad
\dot{\pi}_a(\mathbf{x},t)=-\frac{\delta H}{\delta \phi_a(\mathbf{x},t)}.
\end{equation}

\paragraph{Example: Klein--Gordon Hamiltonian.}
For $\mathcal{L}_{\mathrm{KG}}$:
\begin{equation}
\pi=\dot{\phi},
\end{equation}
and
\begin{equation}
\mathcal{H}
=\frac{1}{2}\pi^2+\frac{1}{2}(\nabla\phi)^2+\frac{1}{2}m^2\phi^2,
\qquad
H=\int d^3x\left[\frac{1}{2}\pi^2+\frac{1}{2}(\nabla\phi)^2+\frac{1}{2}m^2\phi^2\right].
\end{equation}

\subsection{Lorentz Invariance}
A theory is Lorentz invariant if the action is invariant under
\begin{equation}
x^\mu \to x'^\mu=\Lambda^\mu_{\ \nu}x^\nu,
\end{equation}
with fields transforming in some representation:
\begin{equation}
\phi_a(x)\to \phi'_a(x') = D(\Lambda)_{ab}\,\phi_b(x).
\end{equation}
The action is invariant if
\begin{equation}
S[\phi]=S[\phi'].
\end{equation}

For a scalar field ($D(\Lambda)=1$):
\begin{equation}
\phi'(x')=\phi(x).
\end{equation}

\paragraph{Lorentz-covariant building blocks.}
Common invariant terms in $\mathcal{L}$:
\begin{equation}
\phi^2,\qquad \partial_\mu\phi\,\partial^\mu\phi,\qquad
\phi^\ast\phi,\qquad \partial_\mu\phi^\ast\,\partial^\mu\phi.
\end{equation}

\subsection{Spacetime Translations and Energy--Momentum Tensor}
For an infinitesimal translation
\begin{equation}
x^\mu \to x'^\mu = x^\mu + \epsilon^\mu,
\end{equation}
the action is invariant if $\mathcal{L}$ has no explicit $x$ dependence.

The \emph{canonical energy--momentum tensor} is
\begin{equation}
T^{\mu}_{\ \nu}
=
\frac{\partial\mathcal{L}}{\partial(\partial_\mu\phi_a)}\partial_\nu\phi_a
-\delta^\mu_{\ \nu}\mathcal{L}.
\end{equation}
Conservation law:
\begin{equation}
\partial_\mu T^{\mu}_{\ \nu}=0.
\end{equation}

Energy and momentum:
\begin{equation}
P^\nu = \int d^3x\, T^{0\nu},
\qquad
H=P^0=\int d^3x\,T^{00}.
\end{equation}

\paragraph{Remark (symmetrization).}
The canonical $T^{\mu\nu}$ is not always symmetric.
One can construct an improved symmetric tensor (Belinfante):
\begin{equation}
T^{\mu\nu}_{\mathrm{sym}} = T^{\mu\nu} + \partial_\lambda B^{\lambda\mu\nu},
\end{equation}
where $B^{\lambda\mu\nu}$ is antisymmetric in $\lambda\leftrightarrow \mu$,
so $\partial_\mu T^{\mu\nu}_{\mathrm{sym}}=0$ still holds.

\subsection{Noether's Theorem (General Form)}
Suppose the action is invariant under a continuous transformation with
\begin{equation}
\phi_a(x)\to \phi_a(x)+\delta\phi_a(x),
\qquad
x^\mu\to x^\mu+\delta x^\mu,
\end{equation}
such that the Lagrangian changes by a total derivative:
\begin{equation}
\delta\mathcal{L} = \partial_\mu K^\mu.
\end{equation}
Then there exists a conserved current $j^\mu$:
\begin{equation}
\partial_\mu j^\mu=0,
\end{equation}
with conserved charge
\begin{equation}
Q = \int d^3x\, j^0,
\qquad
\frac{dQ}{dt}=0.
\end{equation}

A useful expression for the Noether current is:
\begin{equation}
j^\mu
=
\frac{\partial\mathcal{L}}{\partial(\partial_\mu\phi_a)}\,\delta\phi_a
+\mathcal{L}\,\delta x^\mu
- K^\mu.
\end{equation}

\subsection{Internal Symmetry Example: Global $U(1)$}
Complex scalar field with
\begin{equation}
\mathcal{L}=\partial_\mu\phi^\ast\,\partial^\mu\phi - m^2\phi^\ast\phi.
\end{equation}
Global phase symmetry:
\begin{equation}
\phi \to e^{i\alpha}\phi,
\qquad
\phi^\ast \to e^{-i\alpha}\phi^\ast,
\qquad
\alpha=\text{const}.
\end{equation}
Infinitesimal:
\begin{equation}
\delta\phi = i\alpha\phi,
\qquad
\delta\phi^\ast=-i\alpha\phi^\ast.
\end{equation}
Noether current:
\begin{equation}
j^\mu
=
i\left(\phi^\ast\partial^\mu\phi-(\partial^\mu\phi^\ast)\phi\right).
\end{equation}
Conservation:
\begin{equation}
\partial_\mu j^\mu=0.
\end{equation}
Conserved charge:
\begin{equation}
Q=\int d^3x\, j^0.
\end{equation}

\subsection{Rotations and Angular Momentum (Useful Formulas)}
Lorentz invariance implies conservation of angular momentum (including boosts).
Define the Lorentz current:
\begin{equation}
M^{\lambda\mu\nu}
=
x^\mu T^{\lambda\nu} - x^\nu T^{\lambda\mu}
+\text{(spin part, if field has spin)}.
\end{equation}
Conservation:
\begin{equation}
\partial_\lambda M^{\lambda\mu\nu}=0.
\end{equation}
Angular momentum generators:
\begin{equation}
J^{\mu\nu} = \int d^3x\, M^{0\mu\nu}.
\end{equation}

\subsection{Classical Poisson Brackets for Fields}
At equal time $t$:
\begin{equation}
\{\phi_a(\mathbf{x},t),\pi_b(\mathbf{y},t)\}
=\delta_{ab}\,\delta^{(3)}(\mathbf{x}-\mathbf{y}),
\end{equation}
\begin{equation}
\{\phi_a(\mathbf{x},t),\phi_b(\mathbf{y},t)\}=0,
\qquad
\{\pi_a(\mathbf{x},t),\pi_b(\mathbf{y},t)\}=0.
\end{equation}
Time evolution of any functional $F[\phi,\pi]$:
\begin{equation}
\dot{F}=\{F,H\}+\frac{\partial F}{\partial t}.
\end{equation}

\subsection{Useful Lagrangians (Most Common)}
Real scalar ($\phi^4$ theory):
\begin{equation}
\mathcal{L}
=\frac{1}{2}\partial_\mu\phi\,\partial^\mu\phi
-\frac{1}{2}m^2\phi^2
-\frac{\lambda}{4!}\phi^4.
\end{equation}

Complex scalar with interaction:
\begin{equation}
\mathcal{L}
=\partial_\mu\phi^\ast\,\partial^\mu\phi
-m^2\phi^\ast\phi
-\frac{\lambda}{4}(\phi^\ast\phi)^2.
\end{equation}

\subsection{Dimensional Analysis (Quick Exam Tool)}
In $d=4$ spacetime dimensions:
\begin{equation}
[\mathcal{L}] = 4.
\end{equation}
For a real scalar kinetic term:
\begin{equation}
\mathcal{L}\sim (\partial\phi)^2
\quad\Rightarrow\quad
[\phi]=1.
\end{equation}
Then:
\begin{equation}
[m]=1,\qquad [\lambda_{\phi^4}]=0.
\end{equation}
So $\phi^4$ coupling is dimensionless (renormalizable in $4d$).

\subsection{Summary (What to Remember)}
\begin{itemize}
\item Field equations follow from $\delta S=0$:
\[
\frac{\partial\mathcal{L}}{\partial \phi_a}
-\partial_\mu\left(\frac{\partial\mathcal{L}}{\partial(\partial_\mu\phi_a)}\right)=0.
\]
\item Canonical momentum: $\pi_a=\partial\mathcal{L}/\partial\dot{\phi}_a$.
\item Hamiltonian density: $\mathcal{H}=\pi_a\dot{\phi}_a-\mathcal{L}$.
\item Noether theorem: continuous symmetry $\Rightarrow$ conserved current $\partial_\mu j^\mu=0$.
\item Translation invariance $\Rightarrow$ conserved $T^{\mu\nu}$ and $P^\mu$.
\end{itemize}
